\chapter{Overview}




\section{Intoduction to numerical simulation with TiberCAD{} }

TiberCAD{} is  a multiphysics software tool; it  comprises  a  set  of  solvers, called   simulation models, each one describing a physical problem to be solved, e.g. DriftDiffusion (to solve Poisson and DriftDiffusion equations), EFASchroedinger (to solve Schroedinger equation in envelope function  approximation), Macrostrain (to calculate macroscopical strain with an elastic model) and others.

Similarly to  other device CADs, TiberCAD{} requires that one follows
 a three-step procedure.

 In a first step the device geometry must be sketched, giving all the geometrical information needed by the simulations. This can be performed by means of a   text file or with the help of a graphical tool.
During  this  procedure one or more {\em mesh regions} and {\em boundary  regions} have  to  be  defined: in a following stage the  {\em mesh regions} will  be  associated to  materials and {\em boundary  regions} to device  contacts (boundary  conditions  in  general).
 %The main output of this procedure is a geometry file that specifies one or more {\em mesh regions}. 

The second step consists in running a {\em mesher} tool, which reads the geometry file and sets up the computational mesh used to discretize the partial differential equations representing the physical models to be solved. For this procedure \tibercad  makes use of the GPL software {\em gmsh}. 
Optionally, mesh outout provided by other meshers, such as the one included in  ISE-TCAD tool, are also supported. 
The output of the meshing  procedure is a mesh file that contains information about the space discretization as well as the {\em mesh regions} and the {\em boundary  regions}. 

In the last step the actual simulations are performed. Together with the mesh information (comprised in the  mesh file ), TiberCAD{}  requires an input file which associates {\em materials} to {\em mesh regions}, defines the physical properties and physical models to be applied, and the type of calculations to be performed. 

Details about the modeler and mesher tools can be found in the specific user-manuals. Here we deal primarily with the TiberCAD{}   input file. However, in discussing examples of 1D, 2D and 3D simulations, we will also describe in some detail the geometry input files used to run {\em gmsh}.  
\subsection{ Input file  structure}
%The  execution of  the  requested  simulations  is  controlled  by  a  an  input file ,  e.g.  input.tib 
 TiberCAD{}   input file is  a  text file  which  includes  a  description of  the  device  structure,  the  definition of  the  solvers  to  be  executed,  with  all  the  relevant  physical and  numerical  parameters for  each  of  them.


The input file is organized into several sections, each describing a different aspect of the problem to be solved. The strategy employed is similar to other commercial T-CAD tools and requires some practice to reach a good level of familiarity. We strongly suggest to read first the following chapters (Getting  Started  1 and  2)  in this manual and then study the input files provided in the examples directory and try to  modify them, before writing your own  input file from scratch. The example files touch all current features implemented in the code.

Let's see  an  overview  of  the  main  features of the  input file 
The core of  the input file comprises three sections, called \textbf{Device}, \textbf{Models} and \textbf{Simulation}:

\begin{verbatim}
$Device
{
	Region Si_channel
	{
		
		material= Si	
		doping = 1e16
	}
	Region gate_oxide
	{
		
		material= SiO2
	}
	...
	
}

$Models
{
  model driftdiffusion
  {
    simulation_name= dd
  }
  
  ...
  
}

$Simulation
{
	solve= dd
}
\end{verbatim}

The section \textbf{Device} is used to associate one or more mesh regions to a material and to a set of   physiscal properties such as doping concentrations, doping levels, etc.
 The section
\textbf{Model} is used to define the physical models to be solved. Each physical model can be applied to the whole  device, or to a set of regions of  the  device,  (defined in \textbf{Device}   section). 
Finally, \textbf{Simulation}  section states the type of calculation to  be  executed,  that is  the  \textit{simulation} to  be  \textit{solved}.
These  sections  will  be  described  in full details  in  the  following.

Besides these  main  sections, there  other  other  two,  called  \textbf{Solver} and  \textbf{Physics},  where  some  parameters can  be  set,  respectively  for the  numerical solvers and  for  the  physical  models.
The  aim  of  these  sections is  to  give  the user the  maximum  of  flexibility to  tune his/her  simulation.
Especially  regarding  the  Solver  case,  values  of  the  numerical  parameters  have  been already  tuned for  each  application and   should  be  modified only  by an  advanced  user. 




% One \textbf{simulation} is a particular set of equations, boundary conditions, physical parameters, solver parameters, which describes a physical problem to be solved by \TiberCAD{}.
% A valid \TiberCAD{} \textbf{simulation} must belong to one of the predefined \TiberCAD{} \textbf{simulation} \textbf{models}
% (see section \ref{section:models}).



% To create a \TiberCAD{} \textbf{simulation}, we first have to declare the \TiberCAD{} \textbf{model} class to which our simulation belongs: 
% 
% 
% \begin{verbatim}
% $Models
% {
%   model driftdiffusion
%   {
%     ...
% \end{verbatim}
% 
% Here, we declare that the simulation to be created will belong to the \textbf{model} class driftdiffusion (\textit{model driftdiffusion})
% In the next \textit{Options} block, we define the \textbf{name} of the \TiberCAD{} \textbf{simulation} (\textit{simulationname = }) and the \TiberCAD{} regions to which it will be applied ( e.g., \textit{physical\_regions = all}, where 'all' means the whole device) 
% 
% In general, several \TiberCAD{} \textbf{simulations} belonging to the same \textbf{model} can be created; each of them must have a different name.   As we will see in the following (see \ref{section:solver},\ref{section:physics}) , it is possible to specify physical and solver parameters for a single \TiberCAD{} \textbf{simulation}, by referring to its  \textbf{name}. Anyway, it is also possible to specify parameters \textbf{common} to all the \TiberCAD{} \textbf{simulations} belonging to the same \textbf{model}, by referring to the name of the \TiberCAD{} model instead (in this case \textit{'driftdiffusion' }).




\section{Definition of physical and boundary regions in  TiberCAD{} }

\label{overview_boundaries}

Let'see  now  in  more  detail how  to  associate  \textit{physical}  information to  the  regions  of  the  device  model and  how  to  define  the  boundary  regions,  that  is  contacts or, in general,  regions  where  some  kind  of  boundary  condition is  to be  applied.

When \TiberCAD{}  is  run,  it  reads  the  mesh  file which  contains  the finite element  grid  which meshes the  geometrical description of  the  device or  nanostructure,  and  which will be  the  basis of  PDE  discretization.


As we  have  seen  before, to  execute the  proper  simulations, \TiberCAD{} needs  some  information about  the   physical and boundary regions  associated with the  mesh.
A physical region associates all  the  elements  corresponding to  an  homogeneous  part  of  the  device (usually  related to  the  same  material or  doping).   In  \TiberCAD{}, these regions  are  referred to as \textbf{mesh\_regions}.
\
%\subsection{Definition of contacts: boundary regions in  TiberCAD{} }

As for boundary regions, they  are  needed to  specify   boundary  conditions (b.c.)  for  the  solution of  the PDEs  of  our  simulation.
By  default,  to all   the  external boundary  of  the  device  a  Neumann b.c. is  imposed,   meaning  null  derivative  of  electric  field  and  zero  flux  of  current  normal to  the  boundary.  These are  the  usual  b.c.  applied in  the   simulation  of  electronic  devices;  in   particular, these  conditions are  implicitly  satisfied by using  the  finite  element  formulation..
%
Usually,  however, one  needs  to  impose   also   specifical  b.c.  to  the  device,  relative,  most  often, to  contacts of  some  kind (ohmic, schottky),  but  also   heat and  temperature b.c.  or  reference substrates  for  strain  calculations.
These  regions,  (constituted by  surfaces, lines or  points,  respectively  for  3D,  2D  and  1D  simulations) are  called in  \TiberCAD{}  \textbf{boundary regions}.


It  is  important to  know  that  the  information about  the   physical and boundary regions  must  be  present in  the  mesh  file  before it  is  read  by  \TiberCAD{},  and  thus  have  to  be  produced by  making  use of  the  modeling/mesher  software.
As  for  now,  \TiberCAD{} supports  the  mesh  output  of  the  following   software tools:  \textbf{GMSH v.2} and  \textbf{ISE-TCAD  v.9.5}.

\subsection{Using ISE-TCAD }

By  means  of  the  utility DEVISE    of  ISE-TCAD  v.9.5,  it  is  possible  to  design and  mesh  a  device;  after the  meshing  has  been    successfully performed,  an  output  file is  produced,   with  the  extension  \textit{.grd}.
This  file  contains    the  description of  the  mesh and  also  the  list  of  the  user defined  material  regions    and  contact  regions.
By  reading  this  \textit{.grd} file  in  \TiberCAD{},  one  can  refer  to  the  ISE TCAD material  regions,  simply  with  the  user-defined name,  which is  present  in  the    \textit{.grd}  file.  This  name  should  be  unique in  the  whole  device.
In the  same way,  ISE TCAD Boundary regions (\textit{Contacts})  can  be  referred to in   \TiberCAD{}  by  means  of  their  user-defined name, present in  the   \textit{.grd} output  file,  too.

\subsection{Using GMSH }

If  GMSH  program  is  used  to  model and  mesh  the  device,  a  bit  more  care  has  to  be  taken. 
Here  we   introduce  the  procedure  to  be  followed; in  the  following tutorials  (Getting  started...  the  subject  will  be  considered  in detail  with  a  step-by-step  description.  
See also the GMSH user manual (\url{http://geuz.org/gmsh}) for  further  details. 

In  the  context of  GMSH,  it  is  possible to  define  several 1,  2  and  3D \textit{Physical Entities}.
These \textit{Physical Entities} allow  to  associate one or more  geometrical  entities to  a  single  numerical  ID
or, better, to  a string label,
so  that  several  
\textbf{mesh\_regions}  and  \textbf{boundary regions} can  be  defined and  referred  to  by  TiberCAD{}.  
Here is  a  simple example  of   a  script  to  generate a  1D geometrical  model ( \textit{*.geo})  file in  GMSH: (see also  chapter \ref{chap:Gettingstarted1D})

\begin{verbatim}
Point(1) = {-25,0,0,0.5};
Point(2) = {0,0,0,0.002};
Point(3) = {25,0,0,0.5};
Line(1) = {1,2};
Line(2) = {2,3};


Physical Line("p_side") = {2};
Physical Line("n_side") = {1};
Physical Point("cathode") = {3};
Physical Point("anode") = {1};
\end{verbatim}

Here,  first the   geometrical  entities  \textbf{Point}s  and  \textbf{Line}s  are  defined. 

In the definition of the geometrical \textbf{Points},  the three first expressions inside the braces on the right hand side give the three X, Y and Z coordinates of the point ; the last expression (0.5 or  0.002 in  this  example) sets the \textbf{characteristic mesh length} at that point, that is the \textbf{size} of a mesh element, defined as the length of the segment for a line mesh element, the radius of the circumscribed circle for a triangle mesh element and the radius of the circumscribed sphere for a tetrahedron mesh element.
Thus, the smaller is the value of the \textbf{characteristic mesh length}, the greater is the mesh density close to that point. The size of the mesh elements will
then be computed in GMSH by linearly interpolating these characteristic lengths in the whole mesh.

In the definition of a geometrical \textbf{Line},  the two expressions inside the braces on the right hand side  give the identification numbers of the start and end \textbf{Points} of the line.


 Then,  two  physical  regions are  defined, each  associated to  one  of  the  two  geometrical  entities:  \textbf{Physical Line(n\_side)}  and \textbf{ Physical Line(p\_side)}. 
The expression(s) inside the braces on the right hand side  give, in general, the identification numbers of all the geometrical lines that need to be grouped inside the \textbf{Physical Line}.

 In this  way,  these  physical  regions are made  available for  \TiberCAD{},  and  will be  used  to  associate them  to  a  \TiberCAD{} region through the  keyword \textbf{Region} ,  as  follows:
\begin{verbatim}
Region n_side
   {
     ............

............

Region  p_side 
   {
     .............

...........
\end{verbatim}


It  is also possible  to  group more than  one  physical region in a  single  Device  Region,  with the  keyword \textbf{mesh\_regions},  as  follows:


\begin{verbatim}
Region reg_1
   {
     mesh_regions  = (region1, region2)

............

Region  reg_2
   {
     mesh_regions  = (region3, region4)

...........
\end{verbatim}

Note  that, in  this case, \textit{region1, region2, region3, region4}  are  the  labels  previously  defined inside  GMSH,  while  \textit{reg\_1 and  reg\_2} are  names chosen  by  user for  his  convenience.


Then,  in  the  GMSH  script,  two  \textbf{Physical Point }  are  defined, \textbf{anode}  and \textbf{ cathode},   and  associated to  the  first  and  to  the   last  point  of  our  1D  device.  These  points  are  needed  to  impose  some  boundary  conditions   and  in this  way  they  are made  available for  TiberCAD{},  and  will be  used  to  associate each  of  them  to  a  boundary  condition  region, through the  keyword 
\textbf{BC\_Region}:

\begin{verbatim}
BC_Regions  
    {
      BC_Region cathode
      {
        .........

......................

	BC_Region anode
        {
          ........

....................

\end{verbatim}

As  before,  it  is  possible  group more  boundary regions in  a  single \textbf{BC\_Region}, with  the  keyword 
 \textbf{BC\_reg\_numb}

\begin{verbatim}
BC_Regions  
    {
      BC_Region cathode
      {
        BC_reg_numb = (contact1,  contact2)

......................

	
....................

\end{verbatim}

Again,  in  this second  case,  \textit{cathode}  is  a  label  chosen  by user,  while  \textit{contact1} and  \textit{contact2} are    labels  previously  defined inside  GMSH.


In  2D  case,  a  set  of  \textbf{Physical Surface}  will  be  defined  to  be  used  as   \textbf{mesh\_regions},  while  
\textbf{Physical Line}  is  used  for  \textbf{boundary regions}.

Finally,  in  3D  case,  \textbf{Physical Volume} is   used  to  define  \textbf{mesh\_regions},  while \textbf{Physical Surface} is  used  to  define  \textbf{boundary regions}.




\section{Simulation  environments }

\TiberCAD{} allows  to  compute  different  physical models in  different parts of  a  device or  nanostructure by  coupling in a general way different  \textit{simulation  environments}.
 A \textit{simulation  environment} is  composed  by  all  the  physical regions to  which  a  particular model  is  assigned.
A \textit{simulation  environment} is  therefore  defined by  the  mesh  elements belonging to its physical regions and  by  the  \textbf{simulation model} which has  been  associated to  these  regions.
This  association is  made possible  by  the  definition of  \TiberCAD{} \textbf{Regions} and  \textbf{Clusters}.
Different   \textit{simulation  environments} can  have  physical regions  in  common.
In this  way,  each  simulation  is  run  on  a  subset  of  the  device  and  can  be possibly  coupled  (even self-consistently)   with  a simulation  run on  a  different subset  of  the  device corresponding to a \textit{simulation  environment} with a  non-void  intersection with  the  first  one.  The  coupling  of  the  two  simulations is  performed 
%at  the  boundary between  the  two  \textit{simulation  environments}, 
 by  means  of  appropriate  Boundary  Conditions (e.g. Current Density, Voltage, ...).
 In  principle, the  two \textit{simulation  environments}  can    refer to  two simulations at  different scale, e.g. atomistic  tight-binding  and  macroscopic drift-diffusion. This  allows an effective  \textbf{multi-scale} simulation of the  device to be  studied.














